\documentclass[11pt]{article}
\setlength{\parindent}{0pt}
\usepackage{xltxtra}
\usepackage{hyperref}
\hypersetup{hidelinks}
\usepackage{url}
\urlstyle{tt}
\usepackage{xcolor}
\definecolor{CVBlue}{RGB}{23,110,191}
\usepackage{calc}
\usepackage{graphicx}
\usepackage{tikz}
\usetikzlibrary{calc}
\usepackage{fontspec}
\usepackage{xeCJK}
\usepackage{enumitem}
\CJKsetecglue{} %% 取消中文与数字之间的间隙
%% 主文档字体设置
\setmainfont[
    Path = fonts/Main/,
    Extension = .otf,
    BoldFont = texgyretermes-bold.otf, % 加粗字体
]{texgyretermes-regular.otf} % 正文字体
% 中文字体设置
\setCJKmainfont[
    Path = fonts/hansans/,
    Extension = .ttf,
    BoldFont = NotoSansSC-Bold.ttf, % 加粗字体
]{NotoSansSC-Regular.otf} % 正文字体
\usepackage{relsize}
\usepackage{xspace}
% 使用 fontawesome（部分图标）
\usepackage{fontawesome} 
% A4纸，上下左右边距
\usepackage[
    a4paper,
    left=1.2cm,
    right=1.2cm,
    top=1.5cm,
    bottom=1cm,
    nohead
]{geometry}
\renewcommand{\baselinestretch}{1.5} % 行间距设为1.5
\usepackage{titlesec}
\usepackage{enumitem}
\setlist{noitemsep} % 取消列表项间的额外间距
%\setlist{nosep} % 取消所有垂直间距
\setlist[itemize]{topsep=0.25em, leftmargin=*}
\setlist[enumerate]{topsep=0.25em, leftmargin=*}
% --- 用于控制【不同项目之间】的垂直距离 ---
\newlength{\interProjectSpacing}
\setlength{\interProjectSpacing}{0.9em} % <--- 在此调整项目之间的距离
\newcommand{\projectsep}{\vspace{\interProjectSpacing}}
% --- 用于控制【项目标题】与下方【项目描述】的距离 ---
\newlength{\intraProjectTitleSep}
\setlength{\intraProjectTitleSep}{0.4em} % <--- 在此调整标题和描述的距离
\newcommand{\titlebreak}{\\[\intraProjectTitleSep]}
% --- 用于控制【项目描述】与下方【要点列表】的距离 ---
\newlength{\intraProjectListTopSep}
\setlength{\intraProjectListTopSep}{0.2em} % <--- 在此调整描述和列表的距离
% =======================================================================
\titleformat{\section}         % 定制 \section 命令 
{\large\bfseries\raggedright} % 将 section 标题设置为大号、粗体且左对齐
{}{0em}                      % 可用于为所有 section 添加前缀（如“章节...”）
{}                           % 可用于在标题前插入代码
[{\color{CVBlue}\titlerule}]  % 在标题后插入一条横线
\titlespacing*{\section}{0cm}{*1.6}{*1.2}
\begin{document}
\pagenumbering{gobble}
%%%% 利用tikz来定位照片
\begin{tikzpicture}[remember picture, overlay] 
    \node[anchor = north east] at ($(current page.north east)+(-2cm,-1.2cm)$) {\includegraphics[height=3cm]{avatar.jpg}};
  \end{tikzpicture}%
  %%%% 利用tikz来定位学校Logo，这里只在第一页显示
  \begin{tikzpicture}[remember picture, overlay] 
    \node[anchor = north west] at ($(current page.north west)+(0.5cm,+1.0cm)$) {\includegraphics[height=6cm]{zju.png}};
  \end{tikzpicture}%
\centerline{\LARGE\bfseries{孙梓豪}} 
\centerline{\normalsize{\faPhone\ 111-1111-1111 \quad \faEnvelopeO\ \href{mailto:3230100000@zju.edu.cn}{3230100000@zju.edu.cn}}} 
\centerline{\normalsize{\faGithubSquare\ \href{https://github.com/maksymilan}{https://github.com/maksymilan} \quad \faRssSquare\ \href{https://maksymilan.github.io/}{https://maksymilan.github.io}}} 
    
\section{\makebox[\widthof{\faGraduationCap}][c]{\color{CVBlue}\faGraduationCap}\ 教育背景}    
\textbf{浙江大学} \hfill 2020.9 -- 2024.6\\[0.5em] 
数学与应用数学专业（强基计划）\quad 本科
\projectsep
\textbf{浙江大学} \hfill 2024.9 -- 至今\\[0.5em]
数学专业 \quad 硕士

\section{\makebox[\widthof{\faUsers}][c]{\color{CVBlue}\faUsers}\ 学生工作与社会实践}
% --- 第一个职位 ---
% 将标题行末尾的 \\ 替换为 \titlebreak 命令
\textbf{计算数学研究生党支部副书记} \hfill 2024.09 -- 至今 \titlebreak
职务描述：协助书记负责支部日常建设与党员管理，强化支部政治功能，提升基层党组织的凝聚力与战斗力。
% 在 itemize 的选项中，使用 topsep=\intraProjectListTopSep 来控制上边距
\begin{itemize}[nosep, topsep=\intraProjectListTopSep]
    \item \textbf{组织建设与规范管理}：严格落实“三会一课”制度，主持召开支部大会及支委会，负责党员发展全流程工作（从入党申请人培养到预备党员转正），确保档案整理与组织程序零差错、规范化。
    \item \textbf{“党建+专业”创新实践}：发挥计算数学专业优势，探索“党建+数字化”管理模式，优化支部成员信息管理与党费收缴统计工作，将数据处理技术应用于党务工作场景，有效提升基层党务工作效率与精准度。
    \item \textbf{群众联系与服务}：定期开展谈心谈话，及时掌握研究生同学的思想动态与实际困难，发挥党支部在学术科研与生活保障之间的桥梁纽带作用。
\end{itemize}
% 使用 \projectsep 命令来分隔两个职位
\projectsep
% --- 第二个职位 ---
\textbf{数学科学学院研博会主席} \hfill 2025.09 -- 至今 \titlebreak
职务描述：全面主持研博会工作，重点分管职业发展与就业服务板块，统筹大型活动策划与对外联络。
\begin{itemize}[nosep, topsep=\intraProjectListTopSep]
    \item \textbf{就业指导与服务}：协助学院开展就业经验分享会、企业实践研学、职业能力提升等就业活动，累计服务研究生超500人次；参与编纂《职业发展手册》，切实提升学生的职业规划意识与就业竞争力。
    \item \textbf{活动统筹与文稿撰写}：统筹协调研博会4个部门，参与学院活动的调度与应急处置；撰写活动方案、新闻稿及调研报告等各类文书，具备规范的公文写作基础。
\end{itemize}
\section{\makebox[\widthof{\faCogs}][c]{\color{CVBlue}\faCogs}\ 综合素能}
\begin{itemize}[nosep]
    \item \textbf{办公与写作：} 精通 \textbf{LaTeX} 排版，熟练使用 Office (Word, Excel, PPT) 办公软件，具备规范的机关公文写作习惯。
    \item \textbf{技术能力：} 熟练掌握 \textbf{Python} 等编程语言，具备优秀的数据处理与逻辑分析能力。
    \item \textbf{语言能力：} 英语 \textbf{CET-6}，具备良好的阅读与日常交流能力。
\end{itemize}
\section{\makebox[\widthof{\faGraduationCap}][c]{\color{CVBlue}\faList}\ 获奖情况}
\begin{itemize}
    \item 浙江大学“优秀研究生干部” \hfill 2025.12
    \item 浙江大学“五好研究生” \hfill 2025.12
    \item 浙江大学“灵均领航奖学金” \hfill 2026.10
\end{itemize}
    
\section{\makebox[\widthof{\faInfo}][c]{\color{CVBlue}\faInfo}\ 自我评价}
\begin{itemize}[parsep=0.5ex]
    \item \textbf{政治立场：} 坚持正确政治方向，深学笃用党的创新理论，具有强烈的家国情怀与扎根基层、服务人民的坚定信念。
    \item \textbf{实践动机：} 渴望深入基层党政机关，将数学逻辑思维与数字化技术应用于政务实践，在服务奉献中提升岗位认知。
    \item \textbf{工作态度：} 作风严谨，具备较强的执行力与团队协作精神。
\end{itemize}
\end{document}
